\section{Theory}

\subsection{The 2D Wave Equation}

The first algorithm treats the metal plate as a thin membrane. It has negligible thickness and no elasticity. The waves on the plate is modeled as if they are vibrating freely, with no damping. Therefore, the waves on the plate simply satisfies the 2D wave equation \autocite{PDE}. This is the simplest model that can be constructed to simulate the Chladni plates.

\begin{equation}
   \frac{\partial^2u}{\partial t^2} = c^2 \nabla^2u
\end{equation}

The symbol u is a function of t, time, and x,y which are Cartesian coordinates. It represents the vertical displacement of the wave at the given position and time. The $\nabla^2$ symbol is the 2D laplacian operator. The quantity $c$ is the wave speed.

When the wave is a standing wave on the plate, we can assume that the time and space sections of u is separable, $u(x,y,t) = F(x,y)T(t)$, with the time dependent part $T(t) = cos(\omega t)$. For convenience, x,y, and t will be casted in dimensionless form. Time t will be multiplied by $\frac{L}{c}$. $x$ and $y$ will all be divided by $L$, and u represents a fraction of its amplitude. These three quantities will be expressed in this dimensionless form later on.

In this way, the wave equation can be rewritten into the time independent form:

\begin{equation}
    \nabla^2 u = -k^2 u
\end{equation}

This is essentially a partial differential equation (PDE) and one is trying to find the eigenfunctions of the laplacian operator. One can obtain $\omega$ by multiplying $k$ with $\frac{c}{L}$.

An important detail for this PDE is it's boundary conditions. Normally when solving this type of equation for a vibrating string (or even the Schrodinger's Wave Equation) a Dirichlet boundary condition is assumed. This means that $u = 0$ when $x = 0$ or $x = 1$ (Same for $y$). This means that the wave must disappear at the edges, i.e. clamped strings or infinite potential. However, for a vibrating plate, the edges are allowed to vibrate freely, often vibrating at a maximum or minimum. Therefore, instead of $u(0) = 0$, it is $\frac{du}{dx}_{x=0} = 0$ (Also the same for x = 1 and for y).

Therefore, using this boundary condition, the general solution of this PDE is this:

\begin{equation}
    u(x,y,t) = U_0 cos(\omega_{ms}t)cos(m\pi x)cos(s\pi y)
\end{equation}

Where $U_0$ is the initial displacement (i.e amplitude) of the wave. $m$ and $s$ are the indices for the different modes. $\omega_{ms}$ is the angular frequency of the wave and is determined by $m$ and $s$ via this relationship: $\omega_{ms} = \sqrt{m^2+s^2}$. Note that a value of $\omega$ can correspond to multiple combinations of $m$ and $s$. This means that in the practical, the overall pattern displayed is going to be a superposition of the different degenerate modes.

Another important feature of this solution is that each mode is orthogonal to each other. This is a key thing to remember when dealing with the discretised version of the laplacian, especially when considering the boundary conditions.

\subsection{Numerically Solving the Wave Equation}
Analytical solutions may be convenient for certain situations, however, not any shape has an analytical solution. Sometimes, numerical methods are required. The finite differences method (FDM) can be used to solve this PDE numerically. Dividing the region by N, $dx$ can be turned into $\Delta x$, with $N\Delta x = 1$. This means that the function $u(x)$ is discretised into $u_i$ where i is an integer ranges from 0 to N. Only one dimension is considered in this situation as it simplifies things a lot. The method to raise this to 2D will be shown later.

\subsubsection{Dirichlet Boundary Condition}
Using the central difference approach, the second derivative of a point $u_i$ can be thought as the rate of change between the derivative of the imaginary point $u_{i+0.5}$ and the imaginary point $u_{i-0.5}$.

\begin{equation}
    \frac{\Delta^2 u_i}{\Delta x^2} = (\frac{\Delta u_{i+0.5}}{\Delta x} - \frac{\Delta u_{i-0.5}}{\Delta x}) \frac{1}{\Delta x} \label{eq:centraldiff}
\end{equation}

As the derivatives of $u_{i+0.5}$ and $u_{i-0.5}$ can be defined like this,

\begin{equation}
\frac{\Delta u_{i+0.5}}{\Delta x} = \frac{u_{i+1}-u_i}{\Delta x}; \frac{\Delta u_{i-0.5}}{\Delta x} = \frac{u_i-u_{i-1}}{\Delta x}
\end{equation}

the laplacian operator can be rewritten like this:

\begin{equation}
    \frac{u_{i+1}-2u_{i}+u_{i-1}}{\Delta x^2} = -\omega^2 u_i \label{eq:dlaplacian}
\end{equation}

Or in a matrix form:

\begin{equation}
    \begin{pmatrix}
        1&-2&1
    \end{pmatrix}
    \begin{pmatrix}
        u_{i+1}\\
        u_i\\
        u_{i-1}\\
    \end{pmatrix}
    = -\Delta x^2 \omega^2 u_i
\end{equation}

Because the different $u_i$ can be grouped into one vector like this:

\begin{equation}
    \mathbf{u} = 
\begin{pmatrix}
    u_1\\
    u_2\\
    \vdots\\
    u_{N-1}
\end{pmatrix}
\end{equation}

the matrix equation can therefore be rewritten again into this form:

\begin{equation}
\begin{pmatrix}
    -2    &     1&     0&     0&\cdots&      0\\
    1     &    -2&     1&     0&      & \vdots\\
    0     &     1&    -2&     1&      & \vdots\\
    0     &     0&     1&\ddots&\ddots&      0\\
    \vdots&      &      &\ddots&\ddots&      1\\
    0     &\cdots&\cdots&     0&     1&     -2\\
\end{pmatrix}
\mathbf{u} = -\Delta x^2 \omega^2 \mathbf{u}
\end{equation}

Notice that $u_0$ and $u_N$ are both neglected in this vector as they are all 0 in this boundary condition.

This matrix equation turns the PDE into a simpler eigenvalue and eigenvector problem. Notice that the coefficient $-\Delta x^2\omega^2$ is simply the eigenvalue and $\mathbf{u}$ is the eigenvector of the laplacian matrix. The different eigenvector solutions in this equation is also orthogonal (where the inner product is the dot product instead of the integral). The meaning is that each eigenvector is perpendicular to each other in their ultra high dimensional space (N-1 dim in this case). This is exactly the same ideas as the orthogonal property existing in the continuous, analytical solutions of the equation.

A computer can easily calculate the eigenvector and eigenvalue of this diagonal matrix given enough time and memory.

\subsubsection{Neumann Boundary Condition}

The matrix shown previously is the discrete Dirichlet Laplacian. The case for free edges is the discrete Neumann Laplacian. As mentioned before, Neumann boundary condition assumes zero flux, i.e. $\frac{du}{dx} = 0$, for $x = 0$ or $x=1$. When discretised, a ghost point (non existent but imagined) $u_{-1}$ or $u_{N+1}$ can be used to represent the derivative.

\begin{equation}
    \frac{\Delta u_0}{\Delta x} = \frac{u_1 - u_{-1}}{2\Delta x} = 0 ;
    \frac{\Delta u_N}{\Delta x} = \frac{u_{N+1} - u_{N}}{2\Delta x} = 0
\end{equation}

This implies that $u_1 = u_{-1}$ and $u_{N} = u_{N+1}$. By applying this to \eqref{eq:dlaplacian}, this means that the second derivative of the boundaries is this:

\begin{equation}
    \frac{u_{1}-2u_{0}+u_{-1}}{\Delta x^2} = \frac{2u_1-2u_0}{\Delta x^2}= -\omega^2 u_0
\end{equation}

\begin{equation}
    \frac{u_{N+1}-2u_{N}+u_{N-1}}{\Delta x^2} = \frac{-2u_{N}+2u_{N-1}}{\Delta x^2}= -\omega^2 u_0
\end{equation}

The change to the original matrix equation is that the first row is $\begin{pmatrix}-2&2&0&\cdots&0\end{pmatrix}$ and the last row is $\begin{pmatrix}0&\cdots&0&2&-2\end{pmatrix}$. The matrix also changes from $(N-1)\times(N-1)$ to $(N+1)\times(N+1)$. The full matrix equation is shown bellow:

\begin{equation}
\begin{pmatrix}
    -2    &     2&     0&     0&\cdots&      0\\
    1     &    -2&     1&     0&      & \vdots\\
    0     &     1&    -2&     1&      & \vdots\\
    0     &     0&     1&\ddots&\ddots&      0\\
    \vdots&      &      &\ddots&\ddots&      1\\
    0     &\cdots&\cdots&     0&     2&     -2\\
\end{pmatrix}
\mathbf{u} = -\Delta x^2 \omega^2 \mathbf{u}
\end{equation}

Notice that the vector $\mathbf{u}$ now starts with $u_0$ and ends with $u_N$. There is a problem for this matrix though. The Laplacian operator is a Hermitian operator \autocite{axler2015linear} which gives it the feature of orthogonal eigenfunctions. The corresponding matrix must also be a Hermitian matrix\autocite{horn2013matrix}. This means that the complex conjugate of the transverse of the matrix must be the same as the original matrix. As the matrix shown above is real, this means that the discrete laplacian must have a reflectional symmetry through its main diagonal. The matrix shown above is clearly asymmetrical through the main diagonal. This is not a huge problem mathematically but does put a lot of computational pressure on the computer when it is trying to solve for the eigenvalues and eigenvectors.

To avoid this problem, forward/backward differences can be used instead of central differences method when calculating the second derivatives for the boundaries. Hence, instead of applying the equation shown in \eqref{eq:centraldiff}, one should use this for the left boundary (forward difference): 

\begin{equation}
    \nabla^2 u_0 = (\frac{\Delta u_1}{\Delta x} - \frac{\Delta u_0}{\Delta x})\frac{1}{\Delta x}
\end{equation}

and this for the right boundary (backward difference):

\begin{equation}
    \nabla^2 u_N = (\frac{\Delta u_N}{\Delta x} - \frac{\Delta u_{N-1}}{\Delta x})\frac{1}{\Delta x}
\end{equation}

Using the definition of the Neumann boundary condition, the $\frac{\Delta u_0}{\Delta x}$ term and the $\frac{\Delta u_N}{\Delta x}$ all equals to 0. When backward difference and forward difference is then applied to the left over first derivative terms in the equations respectively, the two second derivatives can be written as follow:

\begin{equation}
    \nabla^2 u_0 = \frac{-u_0 + u_1}{\Delta x^2};
    \nabla^2 u_N = \frac{u_{N-1} - u_N}{\Delta x^2}
\end{equation}

Therefore, the overall discrete Neumann laplacian can be expressed as shown:

\begin{equation}
\begin{pmatrix}
    -1    &     1&     0&     0&\cdots&      0\\
    1     &    -2&     1&     0&      & \vdots\\
    0     &     1&    -2&     1&      & \vdots\\
    0     &     0&     1&\ddots&\ddots&      0\\
    \vdots&      &      &\ddots&    -2&      1\\
    0     &\cdots&\cdots&     0&     1&     -1\\
\end{pmatrix}
\mathbf{u} = -\Delta x^2 \omega^2 \mathbf{u}
\end{equation}

In this way, the diagonally symmetrical (Hermitian) nature of the laplacian matrix is preserved which made it much cleaner to calculate for the eigenvectors and eigenvalues.

\subsubsection{Raising to 2D}
The discretised laplacian matrices mentioned previously are all in 1 dimension. This is done intentionally as it is much easier to work with. However, 1D laplacian might be able to model oscillations on a violin string, it is unable to model the 2D waves on a plate. The laplacian must be transformed to its 2D form to model the Chladni setup correctly. 

The expanded form of the 2D laplacian is shown:

\begin{equation}
    \nabla^2 u = \frac{\partial u}{\partial x} + \frac{\partial u}{\partial y}
\end{equation}

By imagining a L by L space and dividing the space into a N by N mesh, each discrete point on the mesh can be represented by this notation $u_{i,j}$ where $i$ and $j$ both ranges from 0 to N. As the space is divided into a square mesh, a quantity $\Delta h$ can be defined so that $\Delta h = \Delta x = \Delta y$. Therefore, using the equations derived previously, the 2 dimensional second derivative can be expressed as following (a classic 5 point stencil):

\begin{equation}
    \frac{u_{i,j+1}+u_{i+1,j}-4u_{i,j}+u_{i-1,j}+u_{i,j-1}}{\Delta h^2} = -\omega^2 u_{i,j}
\end{equation}

Flattening the 2 dimensional array of $u_{i,j}$ into 1 column of vector $\mathbf{u}$, the matrix equation of the Dirichlet boundary condition is this:

\begin{equation}
    M \mathbf{u} = -\Delta x^2\omega^2 \mathbf{u}
\end{equation}

where the vector $\mathbf{u}$ is 

\begin{equation}
    \mathbf{u} = 
\begin{pmatrix}
    u_{1,1}\\
    \vdots\\
    u_{1,N-1}\\
    u_{2,1}\\
    \vdots\\
    u_{2,N-1}\\
    \vdots\\
    u_{N-1,N-1}
\end{pmatrix}
\end{equation}

$M$ is a $(N-1)^2 \times (N-1)^2$ tridiagonal matrix

\begin{equation}
    M =
\begin{pmatrix}
    B     &     I&     0&     0&\cdots&      0\\
    I     &     B&     I&     0&      & \vdots\\
    0     &     I&     B&     I&      & \vdots\\
    0     &     0&     I&\ddots&\ddots&      0\\
    \vdots&      &      &\ddots&     B&      I\\
    0     &\cdots&\cdots&     0&     I&      B\\
\end{pmatrix}
\end{equation}

with $B$ being another tridiagonal matrix of size $(N-1) \times (N-1)$ that looks like this:

\begin{equation}
B =
\begin{pmatrix}
    -4    &     1&     0&     0&\cdots&      0\\
    1     &    -4&     1&     0&      & \vdots\\
    0     &     1&    -4&     1&      & \vdots\\
    0     &     0&     1&\ddots&\ddots&      0\\
    \vdots&      &      &\ddots&    -4&      1\\
    0     &\cdots&\cdots&     0&     1&     -4\\
\end{pmatrix}
\end{equation}

and $I$ is the identity matrix of the same shape with $B$

The 2 dimensional discrete Neumann Laplacian has a similar equation but with $\mathbf{u}$ being 

\begin{equation}
    \mathbf{u} = 
\begin{pmatrix}
    u_{0,0}\\
    u_{0,1}\\
    \vdots\\
    u_{0,N}\\
    u_{1,0}\\
    \vdots\\
    u_{1,N}\\
    \vdots\\
    u_{N,N}
\end{pmatrix}
\end{equation}

$M$ is now a $(N+1)^2 \times (N+1)^2$ tridiagonal matrix that looks like this\autocite{square_plate}

\begin{equation}
    M =
\begin{pmatrix}
    B_1   &     I&     0&     0&\cdots&      0\\
    I     &   B_2&     I&     0&      & \vdots\\
    0     &     I&   B_2&     I&      & \vdots\\
    0     &     0&     I&\ddots&\ddots&      0\\
    \vdots&      &      &\ddots&   B_2&      I\\
    0     &\cdots&\cdots&     0&     I&    B_1\\
\end{pmatrix}
\end{equation}

with $B_1$ and $B_2$ displayed bellow respectively.

\begin{equation}
B_1 =
\begin{pmatrix}
    -2    &     1&     0&     0&\cdots&      0\\
    1     &    -3&     1&     0&      & \vdots\\
    0     &     1&    -3&     1&      & \vdots\\
    0     &     0&     1&\ddots&\ddots&      0\\
    \vdots&      &      &\ddots&    -3&      1\\
    0     &\cdots&\cdots&     0&     1&     -2\\
\end{pmatrix}
\end{equation}

\begin{equation}
B_2 =
\begin{pmatrix}
    -3    &     1&     0&     0&\cdots&      0\\
    1     &    -4&     1&     0&      & \vdots\\
    0     &     1&    -4&     1&      & \vdots\\
    0     &     0&     1&\ddots&\ddots&      0\\
    \vdots&      &      &\ddots&    -4&      1\\
    0     &\cdots&\cdots&     0&     1&     -3\\
\end{pmatrix}
\end{equation}

These matrices might look horrible at first glance but are actually surprisingly easy to construct. Defining the original 1 dimensional matrix as $D$. Matrix $M$ is simply the kronecker sum \autocite{vanloan2000kronecker} of of 2 $D$s, with the respective boundary conditions.

\begin{equation}
    M = D \oplus D
\end{equation}

The symbol $\oplus$ denotes the kronecker sum and can be expressed as this

\begin{equation}
    D \oplus D = D \otimes I + I \otimes D
\end{equation}

where the symbol $\otimes$ is the kronecker product of a matrix. The exact detail and proof of this will not be shown in here.